% Chtholly Tree (Old Driver Tree)

An efficient data structure for processing range assignment queries on arrays using \lstinline|std::set|. It maintains segments of identical values as elements. Average time complexity is $O(M \log \log N)$ for random inputs where segment merge/assign operations occur frequently.

\begin{lstlisting}
struct Node {
    int l, r;
    mutable long long val;
    Node(int l, int r = -1, long long val = 0) : l(l), r(r), val(val) {}
    bool operator<(const Node& o) const { return l < o.l; }
};

struct ChthollyTree {
    set<Node> odt;

    void init(const vector<long long>& a) {
        odt.clear();
        for (int i = 1; i < a.size(); i++) {
            odt.insert(Node(i, i, a[i]));
        }
    }

    // Splits the segment containing pos into [l, pos-1] and [pos, r]
    // Returns an iterator pointing to the newly created segment starting at pos
    auto split(int pos) {
        auto it = odt.lower_bound(Node(pos));
        if (it != odt.end() && it->l == pos) return it;
        --it;
        if (it->r < pos) return odt.end();
        int l = it->l, r = it->r;
        long long v = it->val;
        odt.erase(it);
        odt.insert(Node(l, pos - 1, v));
        return odt.insert(Node(pos, r, v)).first;
    }

    // Assigns a single uniform value to the entire range [l, r]
    void assign(int l, int r, long long val) {
        auto itr = split(r + 1);
        auto itl = split(l);
        odt.erase(itl, itr);
        odt.insert(Node(l, r, val));
    }

    // Example range query operation (add value to range)
    void add_range(int l, int r, long long val) {
        auto itr = split(r + 1);
        auto itl = split(l);
        for (; itl != itr; ++itl) {
            itl->val += val;
        }
    }
};
\end{lstlisting}

Always call \lstinline|split(r + 1)| before \lstinline|split(l)| to protect the lower bound iterator validity from being invalidated during structural modification steps.
